\src{252}
tính chất phẳng của các bó thích hợp quả thực bảo đảm rằng các phép toán như
$\operatorname{Tor}_i(\mathcal F,\mathcal F')$ giao hoán với mọi phép đổi cơ
sở. Cùng phương pháp ấy áp dụng cho các phép toán mang tính toàn cục: ảnh trực
tiếp và ảnh trực tiếp dẫn xuất của các bó coherent qua các cấu xạ riêng, các
$\operatorname{Ext}$ toàn cục tương đối với các cấu xạ riêng (xem [5], III,
\S~6), v.v.; mọi phép toán này đều biến các họ bị chặn của những bó thành các
họ bị chặn của những bó. (Lưu ý. -- Ở đây những tiền lược đồ mà trên đó ta lấy
các bó khác nhau có thể thay đổi qua các phép toán thuộc kiểu đang xét.)

Hai phát biểu sau được chứng minh về cơ bản bằng cùng kỹ thuật về tính phẳng;
đối với sự phân tích nguyên sơ trên các thớ của một cấu xạ thuộc kiểu hữu hạn,
xem riêng [5], IV.

\result{MỆNH ĐỀ}{1.2} Cho $E$, $E'$ là những tập hợp bị chặn gồm các lớp bó
trên các thớ của $X/S$, giả sử $X$ riêng trên $S$. Khi đó:

\medskip
\noindent (i) Họ các hạt nhân, đối hạt nhân và ảnh của những đồng cấu
$\mathcal F\longrightarrow\mathcal F'$, trong đó lớp của $\mathcal F$ thuộc
$E$ và lớp của $\mathcal F'$ thuộc $E'$, là bị chặn.

\smallskip
\noindent (ii) Họ các bó $\mathcal F''$ là những mở rộng của một $\mathcal F$
bởi một $\mathcal F'$, trong đó lớp của $\mathcal F$ thuộc $E$ và lớp của
$\mathcal F'$ thuộc $E'$, là bị chặn.

Sau một phép đổi cơ sở thích hợp, có thể giả sử $E$ và $E'$ lần lượt được xác
định bởi một bó coherent $\mathcal G$ và một bó coherent $\mathcal G'$ trên
một $X_T/T$, với $T$ thuộc kiểu hữu hạn trên $S$. Hơn nữa, có thể giả sử một
số giả thiết về tính phẳng được thỏa mãn, kéo theo việc tạo thành các bó
$f_{T*}(\operatorname{Hom}_{X_T}(\mathcal G,\mathcal G'))$ và
$\operatorname{Ext}^{1}_{f_T}(\mathcal G,\mathcal G')$ giao hoán với phép đổi
cơ sở theo một cấu xạ tùy ý $T'\longrightarrow T$. Cũng có thể giả sử các bó
coherent nói trên là tự do địa phương trên $T$. Khi đó, gọi $T_0$ và $T_1$
là các phân thớ vectơ trên $T$ mà bó các mầm tiết diện của chúng lần lượt là
các bó nói trên. Khi ấy ta xác định chính tắc một đồng cấu
$\mathcal G_{T_0}\longrightarrow\mathcal G'_{T_0}$ giữa các bó coherent trên
$X_{T_0}$, và một mở rộng
\[
0\longrightarrow\mathcal G'_{T_1}\longrightarrow\mathcal G''
\longrightarrow\mathcal G_{T_1}\longrightarrow0
\]
gồm các bó coherent trên $X_{T_1}$, có một tính chất phổ dụng hiển nhiên. Bó
thứ hai này xác định một họ đại số bao chứa họ đang xét trong (ii). Điều tương
tự cũng đúng với hạt nhân, đối hạt nhân và ảnh của đồng cấu
